
\begin{obeactivity}
\textbf{Aktivitas 9.1: Membandingkan GOST dengan DES}

1. Buatlah tabel perbandingan GOST dan DES berdasarkan: ukuran blok, panjang
   kunci, jumlah putaran, ukuran kotak-S, dan jenis pergeseran yang dipakai.

2. Jelaskan mengapa panjang kunci 256 bit membuat pencarian kunci menyeluruh
   pada GOST lebih sukar daripada pada DES.

3. Jelaskan mengapa keunggulan itu tidak berlaku terhadap kriptanalisis
   struktural. Hubungkan jawaban Anda dengan
   Subbab~\ref{subsec:kriptanalisis-gost}.

4. Sebutkan tiga parameter yang dapat dipilih pada RC5 dan nilai-nilai yang
   dibolehkan untuk setiap parameter.
\end{obeactivity}

\begin{obeactivity}
\textbf{Aktivitas 9.2: Menjalankan Jadwal Kunci RC5 dengan Tangan}

1. Ambil kunci satu byte \code{00} dan tetapkan $w = 32$, $r = 0$, sehingga
   $t = 2$ dan $c = 1$.

2. Hitung $n = 3 \cdot \max(t,c)$, lalu kerjakan langkah pencampuran satu per
   satu menurut Persamaan~\eqref{eq:jadwal-kunci-rc5}. Catat nilai $X$, $Y$,
   $S[0]$, dan $S[1]$ pada setiap langkah.

3. Bandingkan hasil Anda dengan Tabel~\ref{tab:jejak-jadwal-kunci-rc5}.

4. Ulangi langkah 1--3 untuk $r = 8$. Berapa langkah pencampuran yang diperlukan
   sekarang?
\end{obeactivity}

Fungsi pembangkitan kunci putaran pada Aktivitas 9.2 dapat diperiksa melalui
program pada Listing~\ref{lst:rc5-jadwal}. Program itu mengikuti ketiga tahap
pada Subbab~\ref{sec:rc5-rc6}: menyalin kunci ke larik $L$, menginisialisasi
$S$ dari $P_w$ dan $Q_w$, lalu mencampur $L$ dan $S$ sebanyak
$n = 3 \cdot \max(t,c)$ langkah.

\lstinputlisting[
    language=Python,
    firstline=42,
    lastline=82,
    firstnumber=1,
    caption={Pembangkitan kunci putaran RC5: larik $L$, inisialisasi $S$, dan pencampuran},
    label={lst:rc5-jadwal}
]{\codefile{python/rc5_gost_uji.py}}

\begin{obeactivity}
\textbf{Aktivitas 9.3: Fungsi Putaran GOST dengan Tangan}

1. Salin kedelapan kotak-S pada Tabel~\ref{tab:kotak-s-gost} ke dalam tabel
   Anda sendiri. Periksa bahwa setiap baris merupakan permutasi bilangan 0
   sampai 15.

2. Ambil $R = \code{12345678}$ dan $K = \code{61626364}$. Hitung $R + K$ modulo
   $2^{32}$, pecah hasilnya menjadi delapan nibble, lalu substitusikan setiap
   nibble dengan Tabel~\ref{tab:kotak-s-gost}.

3. Gabungkan kedelapan luaran menjadi bilangan 32 bit, geser sirkuler ke kiri
   11 bit, lalu XOR dengan $L = \code{01020304}$. Bandingkan hasil Anda dengan
   Contoh~\ref{ex:gost-round}.

4. Ulangi perhitungan dengan masukan $S_1$ sebesar $0010$. Tunjukkan bahwa
   luarannya 9.
\end{obeactivity}

\lstinputlisting[
    language=Python,
    firstline=151,
    lastline=179,
    firstnumber=1,
    caption={Kedelapan kotak-S GOST dan fungsi putarannya},
    label={lst:gost-putaran}
]{\codefile{python/rc5_gost_uji.py}}

\begin{obeactivity}
\textbf{Aktivitas 9.4: Menghitung Batas Rekeying}

1. Untuk cipher blok berukuran $n$ bit, banyaknya plainteks yang boleh
   dienkripsi dengan satu kunci dibatasi oleh paradoks ulang tahun, yaitu
   $2^{n/2}$ blok. Hitung batas itu untuk $n = 32$, $64$, dan $128$ bit, baik
   dalam satuan blok maupun byte.

2. Bandingkan batas untuk $n = 64$ dengan data yang diolah sistem nyata,
   misalnya peladen web yang mengenkripsi beberapa terabita sehari. Berapa lama
   satu kunci boleh dipakai sebelum batas itu tercapai?

3. Tentukan berapa banyak data yang dituntut setiap serangan pada
   Tabel~\ref{tab:kompleksitas-serangan-gost}. Apakah data itu masih berada di
   dalam batas rekeying?

4. Jelaskan dengan bahasa Anda sendiri mengapa kompleksitas $2^{178}$ tidak
   boleh diterjemahkan menjadi ``GOST belum pecah'' tanpa syarat.
\end{obeactivity}
